\documentclass[ignorenonframetext,]{beamer}
% =====================================================================
% Shared Beamer preamble for all regression lecture slides
%
% Usage in each lecture file:
%   \documentclass[ignorenonframetext,]{beamer}
%   % =====================================================================
% Shared Beamer preamble for all regression lecture slides
%
% Usage in each lecture file:
%   \documentclass[ignorenonframetext,]{beamer}
%   % =====================================================================
% Shared Beamer preamble for all regression lecture slides
%
% Usage in each lecture file:
%   \documentclass[ignorenonframetext,]{beamer}
%   \input{preamble}
%   \title{...}
%   \begin{document}
%   \frame{\titlepage}
%   ...
% =====================================================================

% ---------------------------------------------------------------------
% Theme and captions
% ---------------------------------------------------------------------
\setbeamertemplate{caption}[numbered]
\setbeamertemplate{caption label separator}{: }
\setbeamercolor{caption name}{fg=normal text.fg}
\beamertemplatenavigationsymbolsempty

% ---------------------------------------------------------------------
% Fonts and encoding (pdfLaTeX, XeLaTeX and LuaLaTeX all supported)
% ---------------------------------------------------------------------
\usepackage{lmodern}
\usepackage{amssymb,amsmath}
\usepackage{ifxetex,ifluatex}
\ifnum 0\ifxetex 1\fi\ifluatex 1\fi=0 % if pdftex
\usepackage[T1]{fontenc}
\usepackage[utf8]{inputenc}
\else % if luatex or xelatex
\ifxetex
\usepackage{mathspec}
\else
\usepackage{fontspec}
\fi
\defaultfontfeatures{Ligatures=TeX,Scale=MatchLowercase}
\fi

\usetheme{Madrid}
\usecolortheme{orchid}
\usepackage{beamerthemesplit}

% use upquote if available, for straight quotes in verbatim environments
\IfFileExists{upquote.sty}{\usepackage{upquote}}{}
% use microtype if available
\IfFileExists{microtype.sty}{%
\usepackage{microtype}
\UseMicrotypeSet[protrusion]{basicmath} % disable protrusion for tt fonts
}{}

% ---------------------------------------------------------------------
% Packages
% ---------------------------------------------------------------------
\usepackage{amsthm}
\usepackage{bm}
\usepackage{bbm, dsfont}
\usepackage{extarrows}
\usepackage{kotex}
\usepackage{natbib}
\usepackage{graphicx}
\usepackage{booktabs}
\usepackage{multirow}
\usepackage{makecell}
\usepackage{array}
\usepackage{arydshln}
\usepackage{hhline}
\usepackage{caption}
\usepackage{comment}
\usepackage{needspace}
\usepackage{listings}
\usepackage{tikz}
\usetikzlibrary{arrows,shapes,calc,positioning}

\newif\ifbibliography

% ---------------------------------------------------------------------
% Graphics: scale images so that they do not overflow the page margins
% by default; explicit \includegraphics options still take precedence.
% ---------------------------------------------------------------------
\makeatletter
\def\maxwidth{\ifdim\Gin@nat@width>\linewidth\linewidth\else\Gin@nat@width\fi}
\def\maxheight{\ifdim\Gin@nat@height>\textheight0.8\textheight\else\Gin@nat@height\fi}
\makeatother
\setkeys{Gin}{width=\maxwidth,height=\maxheight,keepaspectratio}

% ---------------------------------------------------------------------
% Colors
% ---------------------------------------------------------------------
\definecolor{azure}{rgb}{0, 0.4, 0.8}
\definecolor{crimson}{rgb}{0.8, 0.4, 0}
\definecolor{gold}{rgb}{0.72, 0.53, 0.04}
\definecolor{ciel}{rgb}{0.39, 0.58, 0.93}
\definecolor{dgreen}{rgb}{0.01, 0.75, 0.24}

% ---------------------------------------------------------------------
% Layout
% ---------------------------------------------------------------------
% Prevent slide breaks in the middle of a paragraph:
\widowpenalties 1 10000
\raggedbottom

\setlength{\parindent}{0pt}
\setlength{\parskip}{6pt plus 2pt minus 1pt}
\setlength{\emergencystretch}{3em}  % prevent overfull lines
\providecommand{\tightlist}{%
\setlength{\itemsep}{0pt}\setlength{\parskip}{0pt}}
\setcounter{secnumdepth}{0}
\captionsetup[table]{belowskip=0pt,aboveskip=0pt}
\captionsetup[figure]{belowskip=0pt,aboveskip=0pt}

% ---------------------------------------------------------------------
% Theorem-like environments
% (Beamer already provides theorem, corollary, lemma, definition,
%  example, ... ; the short names below are used by the lecture files.)
% ---------------------------------------------------------------------
\theoremstyle{plain}
\newtheorem{thm}{Theorem}[section]
\newtheorem{cor}[thm]{Corollary}
\newtheorem{lem}[thm]{Lemma}
\newtheorem{prop}[thm]{Proposition}
\theoremstyle{definition}
\newtheorem{defn}{Definition}[section]
\newtheorem{remk}{Remark}[section]
\newtheorem{exam}{Example}[section]
\newtheorem{prob}{Exercise}[section]
\newtheorem{solu2}{Solution}[section]
\newtheorem{model}{Model}

% ---------------------------------------------------------------------
% Macros
% ---------------------------------------------------------------------
\newcommand{\E}[1]{{\mathbb E}\left[#1\right]}
\newcommand{\Var}[1]{{\mathrm Var}\left(#1\right)}
\newcommand{\var}{\operatorname{Var}}
\newcommand{\dif}{\operatorname{d}}
\newcommand{\parallelsum}{\mathbin{\!/\mkern-5mu/\!}}
\newcommand\tab[1][1cm]{\hspace*{#1}}
\newcommand{\brtexttt}[1]{\texttt{\color{brown}{#1}}}

% ---------------------------------------------------------------------
% Author and date (common to all lectures)
% ---------------------------------------------------------------------
\author[Himchan Jeong]{\large{Himchan Jeong, PhD, FSA}}
\date{}

%   \title{...}
%   \begin{document}
%   \frame{\titlepage}
%   ...
% =====================================================================

% ---------------------------------------------------------------------
% Theme and captions
% ---------------------------------------------------------------------
\setbeamertemplate{caption}[numbered]
\setbeamertemplate{caption label separator}{: }
\setbeamercolor{caption name}{fg=normal text.fg}
\beamertemplatenavigationsymbolsempty

% ---------------------------------------------------------------------
% Fonts and encoding (pdfLaTeX, XeLaTeX and LuaLaTeX all supported)
% ---------------------------------------------------------------------
\usepackage{lmodern}
\usepackage{amssymb,amsmath}
\usepackage{ifxetex,ifluatex}
\ifnum 0\ifxetex 1\fi\ifluatex 1\fi=0 % if pdftex
\usepackage[T1]{fontenc}
\usepackage[utf8]{inputenc}
\else % if luatex or xelatex
\ifxetex
\usepackage{mathspec}
\else
\usepackage{fontspec}
\fi
\defaultfontfeatures{Ligatures=TeX,Scale=MatchLowercase}
\fi

\usetheme{Madrid}
\usecolortheme{orchid}
\usepackage{beamerthemesplit}

% use upquote if available, for straight quotes in verbatim environments
\IfFileExists{upquote.sty}{\usepackage{upquote}}{}
% use microtype if available
\IfFileExists{microtype.sty}{%
\usepackage{microtype}
\UseMicrotypeSet[protrusion]{basicmath} % disable protrusion for tt fonts
}{}

% ---------------------------------------------------------------------
% Packages
% ---------------------------------------------------------------------
\usepackage{amsthm}
\usepackage{bm}
\usepackage{bbm, dsfont}
\usepackage{extarrows}
\usepackage{kotex}
\usepackage{natbib}
\usepackage{graphicx}
\usepackage{booktabs}
\usepackage{multirow}
\usepackage{makecell}
\usepackage{array}
\usepackage{arydshln}
\usepackage{hhline}
\usepackage{caption}
\usepackage{comment}
\usepackage{needspace}
\usepackage{listings}
\usepackage{tikz}
\usetikzlibrary{arrows,shapes,calc,positioning}

\newif\ifbibliography

% ---------------------------------------------------------------------
% Graphics: scale images so that they do not overflow the page margins
% by default; explicit \includegraphics options still take precedence.
% ---------------------------------------------------------------------
\makeatletter
\def\maxwidth{\ifdim\Gin@nat@width>\linewidth\linewidth\else\Gin@nat@width\fi}
\def\maxheight{\ifdim\Gin@nat@height>\textheight0.8\textheight\else\Gin@nat@height\fi}
\makeatother
\setkeys{Gin}{width=\maxwidth,height=\maxheight,keepaspectratio}

% ---------------------------------------------------------------------
% Colors
% ---------------------------------------------------------------------
\definecolor{azure}{rgb}{0, 0.4, 0.8}
\definecolor{crimson}{rgb}{0.8, 0.4, 0}
\definecolor{gold}{rgb}{0.72, 0.53, 0.04}
\definecolor{ciel}{rgb}{0.39, 0.58, 0.93}
\definecolor{dgreen}{rgb}{0.01, 0.75, 0.24}

% ---------------------------------------------------------------------
% Layout
% ---------------------------------------------------------------------
% Prevent slide breaks in the middle of a paragraph:
\widowpenalties 1 10000
\raggedbottom

\setlength{\parindent}{0pt}
\setlength{\parskip}{6pt plus 2pt minus 1pt}
\setlength{\emergencystretch}{3em}  % prevent overfull lines
\providecommand{\tightlist}{%
\setlength{\itemsep}{0pt}\setlength{\parskip}{0pt}}
\setcounter{secnumdepth}{0}
\captionsetup[table]{belowskip=0pt,aboveskip=0pt}
\captionsetup[figure]{belowskip=0pt,aboveskip=0pt}

% ---------------------------------------------------------------------
% Theorem-like environments
% (Beamer already provides theorem, corollary, lemma, definition,
%  example, ... ; the short names below are used by the lecture files.)
% ---------------------------------------------------------------------
\theoremstyle{plain}
\newtheorem{thm}{Theorem}[section]
\newtheorem{cor}[thm]{Corollary}
\newtheorem{lem}[thm]{Lemma}
\newtheorem{prop}[thm]{Proposition}
\theoremstyle{definition}
\newtheorem{defn}{Definition}[section]
\newtheorem{remk}{Remark}[section]
\newtheorem{exam}{Example}[section]
\newtheorem{prob}{Exercise}[section]
\newtheorem{solu2}{Solution}[section]
\newtheorem{model}{Model}

% ---------------------------------------------------------------------
% Macros
% ---------------------------------------------------------------------
\newcommand{\E}[1]{{\mathbb E}\left[#1\right]}
\newcommand{\Var}[1]{{\mathrm Var}\left(#1\right)}
\newcommand{\var}{\operatorname{Var}}
\newcommand{\dif}{\operatorname{d}}
\newcommand{\parallelsum}{\mathbin{\!/\mkern-5mu/\!}}
\newcommand\tab[1][1cm]{\hspace*{#1}}
\newcommand{\brtexttt}[1]{\texttt{\color{brown}{#1}}}

% ---------------------------------------------------------------------
% Author and date (common to all lectures)
% ---------------------------------------------------------------------
\author[Himchan Jeong]{\large{Himchan Jeong, PhD, FSA}}
\date{}

%   \title{...}
%   \begin{document}
%   \frame{\titlepage}
%   ...
% =====================================================================

% ---------------------------------------------------------------------
% Theme and captions
% ---------------------------------------------------------------------
\setbeamertemplate{caption}[numbered]
\setbeamertemplate{caption label separator}{: }
\setbeamercolor{caption name}{fg=normal text.fg}
\beamertemplatenavigationsymbolsempty

% ---------------------------------------------------------------------
% Fonts and encoding (pdfLaTeX, XeLaTeX and LuaLaTeX all supported)
% ---------------------------------------------------------------------
\usepackage{lmodern}
\usepackage{amssymb,amsmath}
\usepackage{ifxetex,ifluatex}
\ifnum 0\ifxetex 1\fi\ifluatex 1\fi=0 % if pdftex
\usepackage[T1]{fontenc}
\usepackage[utf8]{inputenc}
\else % if luatex or xelatex
\ifxetex
\usepackage{mathspec}
\else
\usepackage{fontspec}
\fi
\defaultfontfeatures{Ligatures=TeX,Scale=MatchLowercase}
\fi

\usetheme{Madrid}
\usecolortheme{orchid}
\usepackage{beamerthemesplit}

% use upquote if available, for straight quotes in verbatim environments
\IfFileExists{upquote.sty}{\usepackage{upquote}}{}
% use microtype if available
\IfFileExists{microtype.sty}{%
\usepackage{microtype}
\UseMicrotypeSet[protrusion]{basicmath} % disable protrusion for tt fonts
}{}

% ---------------------------------------------------------------------
% Packages
% ---------------------------------------------------------------------
\usepackage{amsthm}
\usepackage{bm}
\usepackage{bbm, dsfont}
\usepackage{extarrows}
\usepackage{kotex}
\usepackage{natbib}
\usepackage{graphicx}
\usepackage{booktabs}
\usepackage{multirow}
\usepackage{makecell}
\usepackage{array}
\usepackage{arydshln}
\usepackage{hhline}
\usepackage{caption}
\usepackage{comment}
\usepackage{needspace}
\usepackage{listings}
\usepackage{tikz}
\usetikzlibrary{arrows,shapes,calc,positioning}

\newif\ifbibliography

% ---------------------------------------------------------------------
% Graphics: scale images so that they do not overflow the page margins
% by default; explicit \includegraphics options still take precedence.
% ---------------------------------------------------------------------
\makeatletter
\def\maxwidth{\ifdim\Gin@nat@width>\linewidth\linewidth\else\Gin@nat@width\fi}
\def\maxheight{\ifdim\Gin@nat@height>\textheight0.8\textheight\else\Gin@nat@height\fi}
\makeatother
\setkeys{Gin}{width=\maxwidth,height=\maxheight,keepaspectratio}

% ---------------------------------------------------------------------
% Colors
% ---------------------------------------------------------------------
\definecolor{azure}{rgb}{0, 0.4, 0.8}
\definecolor{crimson}{rgb}{0.8, 0.4, 0}
\definecolor{gold}{rgb}{0.72, 0.53, 0.04}
\definecolor{ciel}{rgb}{0.39, 0.58, 0.93}
\definecolor{dgreen}{rgb}{0.01, 0.75, 0.24}

% ---------------------------------------------------------------------
% Layout
% ---------------------------------------------------------------------
% Prevent slide breaks in the middle of a paragraph:
\widowpenalties 1 10000
\raggedbottom

\setlength{\parindent}{0pt}
\setlength{\parskip}{6pt plus 2pt minus 1pt}
\setlength{\emergencystretch}{3em}  % prevent overfull lines
\providecommand{\tightlist}{%
\setlength{\itemsep}{0pt}\setlength{\parskip}{0pt}}
\setcounter{secnumdepth}{0}
\captionsetup[table]{belowskip=0pt,aboveskip=0pt}
\captionsetup[figure]{belowskip=0pt,aboveskip=0pt}

% ---------------------------------------------------------------------
% Theorem-like environments
% (Beamer already provides theorem, corollary, lemma, definition,
%  example, ... ; the short names below are used by the lecture files.)
% ---------------------------------------------------------------------
\theoremstyle{plain}
\newtheorem{thm}{Theorem}[section]
\newtheorem{cor}[thm]{Corollary}
\newtheorem{lem}[thm]{Lemma}
\newtheorem{prop}[thm]{Proposition}
\theoremstyle{definition}
\newtheorem{defn}{Definition}[section]
\newtheorem{remk}{Remark}[section]
\newtheorem{exam}{Example}[section]
\newtheorem{prob}{Exercise}[section]
\newtheorem{solu2}{Solution}[section]
\newtheorem{model}{Model}

% ---------------------------------------------------------------------
% Macros
% ---------------------------------------------------------------------
\newcommand{\E}[1]{{\mathbb E}\left[#1\right]}
\newcommand{\Var}[1]{{\mathrm Var}\left(#1\right)}
\newcommand{\var}{\operatorname{Var}}
\newcommand{\dif}{\operatorname{d}}
\newcommand{\parallelsum}{\mathbin{\!/\mkern-5mu/\!}}
\newcommand\tab[1][1cm]{\hspace*{#1}}
\newcommand{\brtexttt}[1]{\texttt{\color{brown}{#1}}}

% ---------------------------------------------------------------------
% Author and date (common to all lectures)
% ---------------------------------------------------------------------
\author[Himchan Jeong]{\large{Himchan Jeong, PhD, FSA}}
\date{}


\title{Generalized Linear Mixed Models}

\begin{document}
\frame{\titlepage}

\begin{frame}{Table of Contents}
\begin{itemize}
	\item Part 1: Introduction of Generalized Linear Mixed Models
	\vspace{8pt}
	\item Part 2: Credibility premium of $S$ with the \textbf{static} random effects
	\vspace{8pt}
	\item Part 3: Credibility premium of $S$ with the \textbf{dynamic} random effects
\end{itemize}
\end{frame}

%======================================================================
% Part 1
%======================================================================
\section{Introduction of Generalized Linear Mixed Models}

\begin{frame}
\centering
\Large\textbf{Part 1: Introduction of\\ Generalized Linear Mixed Models}
\end{frame}

% ---------------------------------------------------------------------
\begin{frame}[allowframebreaks]
\frametitle{Motivation}

Recall the basic assumptions of the linear models where $y_1,\ldots,y_m$ are the observed response variable and $\mathbf{x}_1,\ldots, \mathbf{x}_m \in \mathbb{R}^{p \times 1}$ are the individual covariates  (here $Z \in \mathbb{R}^{a \times b}$ means that $Z$ is a $(a \times b)$ real-valued matrix):
\begin{enumerate}
	\item $E[y_i] = \mathbf{x}_i' \bm{\beta}$ for $\bm{\beta} = (\beta_1, \ldots, \beta_p) \in \mathbb{R}^{p\times 1}$,
	\item $ \{\mathbf{x}_1, \mathbf{x}_2, \ldots, \mathbf{x}_m \}$ are non-random,
	\item $Var(y_i) = \sigma^2$,
	\item $\{ y_1, y_2, \ldots, y_m  \}$ are independent,
	\item $\{ y_1, y_2, \ldots, y_m  \}$ are normally distributed.
\end{enumerate}
Such situations can be written in short as follows:
\begin{equation} \label{p1:linmod}
y_i = \mathbf{x}_i' \bm{\beta}  + \epsilon_i, \quad \epsilon_i \overset{i.i.d.}\sim \mathcal{N}(0, \sigma^2).
\end{equation}

Note that the first components of $ \{\mathbf{x}_1, \mathbf{x}_2, \ldots, \mathbf{x}_m \}$ are all 1 so that $\beta_1$ corresponds to the \emph{intercept} term in the linear models.

\framebreak

Now, let us assume an example of data with hierarchy so that we observe data points for each of the subjects $i \in \{1,\ldots, m\}$ for time $t=1,\ldots, n_i$. If we adopt the assumptions from the linear models, then we can write them as follows with the modified subscripts:
\begin{enumerate}
	\item $E[y_{i,t}] = \mathbf{x}_{i,t}' \bm{\beta}$ for $\beta \in \mathbb{R}^{p \times 1}$,
	\item $ \{\mathbf{x}_{1,1}, \mathbf{x}_{1,2}, \ldots, \mathbf{x}_{m,n_m} \}$ are non-random,
	\item $Var(y_{i,t}) = \sigma^2$,
	\item {\color{red} $\{ y_{1,1}, y_{1,2}, \ldots, y_{m,n_m}  \}$ are independent,}
	\item {\color{red} $\{  y_{1,1}, y_{1,2}, \ldots, y_{m,n_m}   \}$ are normally distributed.}
\end{enumerate}
What if both of Assumptions 4 and 5 are not valid...? (Note that $\{ y_{i,1}, y_{i,2}, \ldots, y_{i,n_i}  \}$ are from the same subject $i$ over time)

\framebreak

In this regard, one can consider the following representation:
\begin{equation} \label{p1:glmm_intercept}
\mathbb{E}[y_{i,t}|u_i] = \mu_{i,t}, \quad \eta(\mu_{i,t}) = \mathbf{x}_{i,t}'\bm{\beta} + {\color{red}u_i},
\end{equation}
where $u_i$ is a random variable for $i=1,\ldots, m$. Note that it is usual to impose either $\mathbb{E}[u_i]=0$ or $\mathbb{E}[e^{u_i}]=1$.

Although we may incorporate random effects in various ways (not only as the random intercept), for this section we stick to the form in \eqref{p1:glmm_intercept}. Again, $y_{i,t}$ and $y_{i,t'}$ are NOT independent!

\end{frame}

% ---------------------------------------------------------------------
\begin{frame}[allowframebreaks]
\frametitle{Examples}

\begin{example}[Normal GLMM]
Consider the case where $y_{i,t}|u_i$ follows a normal distribution (which belongs to the exponential families) and $\eta(z)=z$ (identity link). In this case, we have
$$
y_{i,t}|u_i \overset{indep}\sim \mathcal{N}(\mathbf{x}_{i,t}'\bm{\beta} + u_i, \sigma^2), \quad \mathbb{E}[u_i]=0,
$$
which is indeed equivalent to $y_{i,t} = \mathbf{x}_{i,t}'\bm{\beta} + u_i + \epsilon_{i,t}$, $\epsilon_{i,t} \overset{i.i.d.}\sim \mathcal{N}(0, \sigma^2)$. Therefore, one can confirm that the GLMMs are natural extensions of the LMMs.
\end{example}

\framebreak

\begin{example}[Bernoulli GLMM - logistic regression]
Consider the case where $y_{i,t}|u_i$ follows a Bernoulli distribution (a special case of the binomial distribution with size 1) and $\eta(z)=\log\left(\frac{z}{1-z}\right)$ (logit link). In this case, we have
$$
y_{i,t}|u_i \overset{indep}\sim \mathcal{B}er\left( \frac{e^{\mathbf{x}_{i,t}'\bm{\beta} + u_i}}{1+e^{\mathbf{x}_{i,t}'\bm{\beta} + u_i}}\right), \quad \mathbb{E}[u_i]=0,
$$
so that $\log\left(\frac{E[y_{i,t}|u_i]}{1-E[y_{i,t}|u_i]}\right) = \mathbf{x}_{i,t}'\bm{\beta} + u_i$, and $\mathbb{P}(y_{i,t}=1|u_i) = \mathbb{E}[y_{i,t}|u_i] = \frac{e^{\mathbf{x}_{i,t}'\bm{\beta} + u_i}}{1+e^{\mathbf{x}_{i,t}'\bm{\beta} + u_i}}$.
\end{example}

\framebreak

\begin{example}[Poisson GLMM]
Consider the case where $y_{i,t}|u_i$ follows a Poisson distribution (which belongs to the exponential families) and $\eta(z)=\log z$ (log link). In this case, we have
$$
y_{i,t}|u_i \overset{indep}\sim \mathcal{P}(e^{\mathbf{x}_{i,t}'\bm{\beta} + u_i}), \quad \mathbb{E}[e^{u_i}]=1,
$$
so that $\log \mathbb{E}[y_{i,t}|u_i] = \mathbf{x}_{i,t}'\bm{\beta} + u_i$ and $\mathbb{E}[y_{i,t}|u_i] = e^{\mathbf{x}_{i,t}'\bm{\beta} + u_i}$.

While it is not strictly required, we often assume that $\Theta_i := e^{u_i} \sim \mathcal{G}(r, 1/r)$ so that $\mathbb{E}[\Theta_i]=1$ and $\mathbb{V}ar(\Theta_i) = 1/r$.
\end{example}

\framebreak

\begin{example}[Gamma GLMM]
Consider the case where $y_{i,t}|u_i$ follows a gamma distribution and $\eta(z)=\log z$ (log link). In this case, we have
$$
y_{i,t}|u_i \overset{indep}\sim \mathcal{G}(1/\psi, \psi \cdot e^{\mathbf{x}_{i,t}'\bm{\beta} + u_i}), \quad \mathbb{E}[e^{u_i}]=1,
$$
so that $\log \mathbb{E}[y_{i,t}|u_i] = \mathbf{x}_{i,t}'\bm{\beta} + u_i$ and $\mathbb{E}[y_{i,t}|u_i] = e^{\mathbf{x}_{i,t}'\bm{\beta} + u_i}$.

While it is not strictly required, we often assume that $\Theta_i := e^{u_i} \sim \mathcal{IG}(r+1, r)$ so that $\mathbb{E}[\Theta_i]=1$ and $\mathbb{V}ar(\Theta_i) = \frac{r}{r-1}$.
\end{example}

\end{frame}

% ---------------------------------------------------------------------
\begin{frame}[allowframebreaks]
\frametitle{Estimation}

\begin{itemize}
	\item Due to the common random effect $u_i$ that affects all observations $y_{i,1}, \ldots, y_{i,n_i}$ simultaneously, $y_{i,1}, \ldots, y_{i,n_i}$ are not marginally independent.
	\item So in the maximum likelihood estimation of $\bm{\beta}$, we need to consider the following joint densities $h(y_{i,1}, \ldots, y_{i,n_i})$, which is given as follows:
	$$
	h(y_{i,1}, \ldots, y_{i,n_i}) = \int \left[ \prod_{t=1}^{n_i} f(y_{i,t}|u_i) \right] \pi(u_i) du_i,
	$$
	where $\pi(u_i)$ is the PDF of $u_i$.
\end{itemize}

\framebreak

\begin{itemize}
	\item Therefore, the loglikelihood is given as follows:
	\begin{align*}
	\ell(\bm{\beta}) &= \sum_{i=1}^m \log h(y_{i,1}, \ldots, y_{i,n_i}; \mathbf{x}_{i,1}, \ldots, \mathbf{x}_{i,n_i} | \bm{\beta}) \\
	&= \sum_{i=1}^m \log \int \left[ \prod_{t=1}^{n_i} f(y_{i,t}|u_i) \right] \pi(u_i) du_i,
	\end{align*}
	and the following score equation
	$$
	S(\bm{\beta}) = \frac{\partial}{\partial \bm{\beta}} \ell(\bm{\beta}) = 0
	$$
	is not easy to be explicitly obtained due to the integrals inside the summation.
\end{itemize}

\framebreak

In general, $\ell(\bm{\beta})$ is not explicitly expressed without the integrals so one can consider the following ways for finding the MLE of $\bm{\beta}$:
\begin{itemize}
	\item EM algorithms: `pretend' that you have complete observations of $u_i$, and iterate to obtain the $\hat{\bm{\beta}}$,
	\item The adaptive Gauss-Hermite approximation: approximate the integral as a finite sum and then evaluate the likelihood.
\end{itemize}
If we assume $u_i$ follows a normal distribution, then we can use the adaptive Gauss-Hermite approximation to find $\hat{\bm{\beta}}$ in \texttt{R} using a package (e.g., \texttt{glmer}).

\framebreak

For the remainder of this chapter, we focus on some special cases where we can obtain an explicit form of $\ell(\bm{\beta})$ with special combinations of $y_{i,t}|u_i$ and $u_i$ (or $y_{i,t}|\Theta_i$ and $\Theta_i$).
\begin{itemize}
	\item Normal-Normal GLMM:
	\begin{align*}
	y_{i,t}|u_i &\overset{indep}\sim \mathcal{N}(\mathbf{x}_{i,t}'\bm{\beta} + u_i, \sigma^2) \text{ for } t=1,\ldots, n_i, \\
	u_i &\overset{indep}\sim \mathcal{N}(0, \phi_u) \text{ for } i=1,\ldots, m. \\
	\Longrightarrow & \ \bm{y}_i \sim \mathcal{N}\left( \mathbf{x}_i \bm{\beta}, \, \phi_u \cdot J_{n_i} + \sigma^2 \cdot I_{n_i} \right),
	\end{align*}
	where $\bm{y}_i = (y_{i,1}, \ldots, y_{i,n_i})' \in \mathbb{R}^{n_i \times 1}$.
\end{itemize}

\framebreak

\begin{itemize}
	\item Poisson-gamma GLMM:
	\begin{align*}
	y_{i,t}|\Theta_i &\overset{indep}\sim \mathcal{P}(\Theta_i \cdot e^{\mathbf{x}_{i,t}'\bm{\beta}}) \text{ for } t=1,\ldots, n_i, \\
	\Theta_i &\overset{indep}\sim \mathcal{G}(r, 1/r) \text{ for } i=1,\ldots, m. \\
	\Longrightarrow & \ \bm{y}_i \text{ follows a \emph{multivariate negative binomial} distribution,}
	\end{align*}
	where $\bm{y}_i = (y_{i,1}, \ldots, y_{i,n_i})' \in \mathbb{R}^{n_i \times 1}$.
\end{itemize}

\framebreak

\begin{itemize}
	\item Gamma-inverse gamma GLMM:
	\begin{align*}
	y_{i,t}|\Theta_i &\overset{indep}\sim \mathcal{G}(1/\psi, \psi \cdot \Theta_i \cdot e^{\mathbf{x}_{i,t}'\bm{\beta}}) \text{ for } t=1,\ldots, n_i, \\
	\Theta_i &\overset{indep}\sim \mathcal{IG}(r+1, r) \text{ for } i=1,\ldots, m. \\
	\Longrightarrow & \ \bm{y}_i \text{ follows a \emph{multivariate generalized Pareto} distribution,}
	\end{align*}
\end{itemize}
where $\bm{y}_i = (y_{i,1}, \ldots, y_{i,n_i})' \in \mathbb{R}^{n_i \times 1}$.

\end{frame}

% ---------------------------------------------------------------------
\begin{frame}[allowframebreaks]
\frametitle{Poisson-gamma GLMM}

\begin{theorem}[Fundamental Theorem of Calculus]
Suppose $\frac{d}{dt}F(t) = f(t)$, then
$$
\int_a^b f(t) dt = F(b) - F(a).
$$
\end{theorem}

\framebreak

\begin{theorem}[Fundamental Theorem of the Loss Models(?)]
If $\Gamma(r)$ is defined in the following way,
$$
\Gamma(r) = \int_0^\infty x^{r-1} e^{-x} dx,
$$
then $\Gamma(r) = (r-1)\Gamma(r-1)$.\\
Especially, if $r$ is a natural number, then $\Gamma(r) = (r-1)!$.
\end{theorem}

\framebreak

Now, we are ready to evaluate the joint PMF of $(y_{i,1}, \ldots, y_{i,n_i})$ in the case of Poisson-gamma GLMM, which is given as follows:
$$
h(y_{i,1}, \ldots, y_{i,n_i}) = \int \left[ \prod_{t=1}^{n_i} f(y_{i,t}|\theta_i) \right] \pi(\theta_i) d\theta_i,
$$
Recall that in the case of Poisson-gamma GLMM, we have
\begin{align*}
f(y_{i,t}|\theta_i) &= \frac{(\theta_i \nu_{i,t})^{y_{i,t}} e^{-\theta_i \nu_{i,t}}}{y_{i,t}!} \text{ for } t=1,\ldots, n_i, \\
\pi(\theta_i) &= \frac{r^r}{\Gamma(r)} \theta_i^{r-1} e^{-r\theta_i},
\end{align*}
where $\nu_{i,t} = \exp\left(\mathbf{x}_{i,t}'\bm{\beta}\right)$.\\
Can you see that $\mathbb{E}[\Theta_i] = \int_0^\infty \theta_i \pi(\theta_i) d\theta_i = 1$?

\end{frame}

% ---------------------------------------------------------------------
\begin{frame}[t]
\frametitle{Multivariate negative binomial distribution}
\footnotesize

Let us first define $\nu_{i,\bullet} = \sum_{t=1}^{n_i} \nu_{i,t}$ and $y_{i,\bullet} = \sum_{t=1}^{n_i} y_{i,t}$.
\setlength{\abovedisplayskip}{4pt}\setlength{\jot}{2pt}
\begin{align*}
& h(y_{i,1}, \ldots, y_{i,n_i}) \\
&= {\color{blue}\int_0^\infty} \left[ \prod_{t=1}^{n_i} f(y_{i,t}|{\color{blue}\theta_i}) \right] {\color{blue}\pi(\theta_i) d\theta_i} \\
& \onslide<2->{= {\color{blue}\int_0^\infty} \left[ \prod_{t=1}^{n_i} \frac{({\color{blue}\theta_i} \nu_{i,t})^{y_{i,t}} e^{-{\color{blue}\theta_i} \nu_{i,t}}}{y_{i,t}!} \right] \frac{r^r}{\Gamma(r)} {\color{blue}\theta_i^{r-1} e^{-r\theta_i} d\theta_i},} \\
& \onslide<3->{= \left[ \prod_{t=1}^{n_i} \frac{(\nu_{i,t})^{y_{i,t}}}{y_{i,t}!} \right] \frac{r^r}{\Gamma(r)} \times {\color{blue}\int_0^\infty \left[ \prod_{t=1}^{n_i} (\theta_i)^{y_{i,t}} e^{-\theta_i \nu_{i,t}} \right] \theta_i^{r-1} e^{-r\theta_i} d\theta_i},} \\
& \onslide<4->{= \left[ \prod_{t=1}^{n_i} \frac{(\nu_{i,t})^{y_{i,t}}}{y_{i,t}!} \right] \frac{r^r}{\Gamma(r)} \times {\color{blue}\int_0^\infty \theta_i^{\sum_{t=1}^{n_i} y_{i,t} + r - 1} e^{-\theta_i (\sum_{t=1}^{n_i} \nu_{i,t} + r)} d\theta_i},} \\
& \onslide<5->{= \left[ \prod_{t=1}^{n_i} \frac{(\nu_{i,t})^{y_{i,t}}}{y_{i,t}!} \right] \frac{r^r}{\Gamma(r)} \times {\color{blue}\int_0^\infty \theta_i^{y_{i,\bullet} + r - 1} e^{-\theta_i (\nu_{i,\bullet} + r)} d\theta_i},}
\end{align*}

\end{frame}

% ---------------------------------------------------------------------
\begin{frame}
\frametitle{Multivariate negative binomial distribution}

By letting $u = \theta_i(\nu_{i,\bullet} + r)$, we have
\begin{align*}
&= \left[ \prod_{t=1}^{n_i} \frac{(\nu_{i,t})^{y_{i,t}}}{y_{i,t}!} \right] \frac{r^r}{\Gamma(r)} \times {\color{blue}\int_0^\infty \theta_i^{y_{i,\bullet} + r - 1} e^{-\theta_i (\nu_{i,\bullet} + r)} d\theta_i}, \\
& \onslide<2->{= \left[ \prod_{t=1}^{n_i} \frac{(\nu_{i,t})^{y_{i,t}}}{y_{i,t}!} \right] \frac{r^r}{\Gamma(r)} \times {\color{blue}\frac{1}{(\nu_{i,\bullet} + r)^{y_{i,\bullet}+r}} \int_0^\infty u^{y_{i,\bullet} + r - 1} e^{-u} du},} \\
& \onslide<3->{= \left[ \prod_{t=1}^{n_i} \frac{(\nu_{i,t})^{y_{i,t}}}{y_{i,t}!} \right] \frac{r^r}{\Gamma(r)} \times {\color{blue}\frac{\Gamma(y_{i,\bullet}+r)}{(\nu_{i,\bullet} + r)^{y_{i,\bullet}+r}}}!}
\end{align*}

\onslide<3->{After such a tedious calculation, we finally removed the integral from the joint PMF $h(y_{i,1}, \ldots, y_{i,n_i})$ and now ready to (eventually) write out $\ell(\bm{\beta})$ explicitly.}

\end{frame}

% ---------------------------------------------------------------------
\begin{frame}[t]
\frametitle{Poisson-gamma GLMM full loglikelihood}
\footnotesize

\vspace{-8pt}
\setlength{\jot}{1pt}
\begin{align*}
\ell(\bm{\beta}) &= \sum_{i=1}^m \log h(y_{i,1}, \ldots, y_{i,n_i}; \mathbf{x}_{i,1}, \ldots, \mathbf{x}_{i,n_i} | \bm{\beta}) \\
& \onslide<2->{= \sum_{i=1}^m \log \int \left[ \prod_{t=1}^{n_i} f(y_{i,t}|\theta_i) \right] \pi(\theta_i) d\theta_i} \\
& \onslide<3->{= \sum_{i=1}^m \left[ \sum_{t=1}^{n_i} (y_{i,t} \log \nu_{i,t} - \log y_{i,t}!) + r \log r \right.} \\
& \onslide<4->{\left. - \log \Gamma(r) + \log \Gamma(y_{i,\bullet}+r) - (y_{i,\bullet}+r)\log(\nu_{i,\bullet}+r) \right],} \\
& \onslide<5->{{\color{purple} = \sum_{i=1}^m \left[ \sum_{t=1}^{n_i} (y_{i,t} \log \nu_{i,t}) - (y_{i,\bullet}+r)\log(\nu_{i,\bullet}+r) \right]}} \\
& \onslide<6->{+ \sum_{i=1}^m \left[ \sum_{t=1}^{n_i} (-\log y_{i,t}!) + \log \Gamma(y_{i,\bullet}+r) \right] + mr\log r - m \log \Gamma(r)}
\end{align*}

\onslide<6->{Note that only {\color{purple}purple} part is relevant for the estimation of $\bm{\beta}$.}

\end{frame}

% ---------------------------------------------------------------------
\begin{frame}[fragile]
\frametitle{Poisson-gamma GLMM full loglikelihood - in \texttt{R}}

\begin{lstlisting}[language=R]
"negll.MVNB" <- function(parm) {

## declarations of variables
p      <- ncol(x);
beta   <- parm[1:p]
r      <- parm[p+1]
nu     <- exp(x %*% beta+log(offset));
dat    <- cbind(id,nu,y);
y_dot  <- aggregate(y~id,dat,sum)[,2]
nu_dot <- aggregate(nu~id,dat,sum)[,2]
m      <- length(unique(id))
\end{lstlisting}

\end{frame}

\begin{frame}[fragile]
\frametitle{Poisson-gamma GLMM full loglikelihood - in \texttt{R}}

\begin{lstlisting}[language=R]
## full loglikelihood
temp1 = sum(y*log(nu))
                - sum( (y_dot+r)* log(nu_dot+r) )
temp2 = -sum(lgamma(y+1)) + sum(lgamma(y_dot+r))
                + m*r*log(r)-m*lgamma(r)
result = -temp1-temp2
return(result)
}
\end{lstlisting}

\end{frame}

\begin{frame}[fragile]
\frametitle{Poisson-gamma GLMM full loglikelihood - in \texttt{R}}

\begin{lstlisting}[language=R]
## optimization in R
init.est <- as.vector(c(init.beta, init.r))
# init.est could be borrowed from marginal models

fit.MVNB <- optim(init.est, negll.MVNB,
                                 hessian=TRUE)

beta.hat <- fit.MVNB$par[1:p]
r.hat <- fit.MVNB$par[p+1]
# r.hat should be positive!
parm.stderr <- sqrt(diag(solve(fit.MVNB$hessian)))
# parm.stderr should be positive!
loglik.MVNB <- -fit.MVNB$value
\end{lstlisting}

\end{frame}

% ---------------------------------------------------------------------
\begin{frame}
\frametitle{Poisson-gamma GLMM - simulation study}

To confirm that the proposed Poisson-gamma GLMM properly works, we generate $\{y_{i,t}\}_{i=1,\ldots,50000, \, t=1,\ldots,6}$ as follows:
\begin{gather*}
y_{i,t}|\Theta_i \overset{indep.}\sim \mathcal{P}(\Theta_i \nu_{i,t}), \ \Theta_i \overset{i.i.d.}\sim \mathcal{G}(r, 1/r), \\
x_{1,i,t} \sim \mathcal{N}(0,1), \ x_{2,i,t} \sim \mathcal{BN}(0.6),
\end{gather*}
and
\begin{align*}
\nu_{i,t} &= m_{i,t} \cdot e^{\mathbf{x}_{i,t}'\bm{\beta}} = \exp(\beta_0 + \beta_1 x_{1,i,t} + \beta_2 x_{2,i,t} + \log m_{i,t}), \\
\beta_0 &= -2.5, \, \beta_1 = 1.5, \, \beta_2 = 0.20, \, r = 5.
\end{align*}
We use the data points from $t=1,\ldots,5$ for the estimation purpose, and those from $t=6$ for the out-of-sample validation purpose, respectively.

\end{frame}

% ---------------------------------------------------------------------
\begin{frame}
\frametitle{Poisson-gamma GLMM - estimation results}

\begin{table}[ht]
\centering
\begin{tabular}{lrrrr}
\toprule
 & \multicolumn{2}{c}{Poisson} & \multicolumn{2}{c}{MVNB} \\
\cmidrule(lr){2-3} \cmidrule(lr){4-5}
 & Estimate & Std. Error & Estimate & Std. Error \\
\midrule
$\beta_0$ & -2.5063 & 0.0110 & -2.5053 & 0.0113 \\
$\beta_1$ & 1.4963 & 0.0096 & 1.4959 & 0.0100 \\
$\beta_2$ & 0.2115 & 0.0120 & 0.2104 & 0.0123 \\
$r$ & - & - & 4.7422 & 0.2686 \\
Loglikelihood & \multicolumn{2}{c}{-85028.36} & \multicolumn{2}{c}{-84811.28} \\
AIC & \multicolumn{2}{c}{170062.71} & \multicolumn{2}{c}{169630.56} \\
BIC & \multicolumn{2}{c}{170094.00} & \multicolumn{2}{c}{169647.42} \\
\hline\hline
\end{tabular}
\end{table}

\end{frame}

\begin{frame}
\frametitle{Poisson-gamma GLMM - estimation results}

We observe that
\begin{itemize}
	\item Interestingly, both the Poisson GLM and Poisson-gamma GLMM (MVNB) produce reliable results in the estimation of $\bm{\beta}$.
	\item In fact, as long as the marginal mean is well-specified (here $\mathbb{E}[y_{i,t}] = \nu_{i,t}$ for both models), then $\hat{\bm{\beta}}$ is consistent (very close to the true $\bm{\beta}$ in case the sample size is very large). Such idea was extensively explored in the statistics literature, which is known as GEE (\href{https://en.wikipedia.org/wiki/Generalized_estimating_equation}{\color{cyan}Generalized estimating equation}, not \href{https://www.google.com/search?q=gee+snsd}{\color{cyan}this!}).
	\item However, it is confirmed that the Poisson-gamma GLMM (MVNB) improves the goodness-of-fit significantly.
\end{itemize}
We will later analyze how these two models are working in terms of the out-of-sample validation (prediction).

\end{frame}

% ---------------------------------------------------------------------
\begin{frame}
\frametitle{Poisson-gamma GLMM - prediction}

Recall that we can obtain the joint probability function of $(y_{i,1}, \ldots, y_{i,n_i})$ as follows:
$$
h(y_{i,1}, \ldots, y_{i,n_i}) = \left[ \prod_{t=1}^{n_i} \frac{(\nu_{i,t})^{y_{i,t}}}{y_{i,t}!} \right] \frac{r^r}{\Gamma(r)} \times {\color{blue}\frac{\Gamma(y_{i,\bullet}+r)}{(\nu_{i,\bullet} + r)^{y_{i,\bullet}+r}}}.
$$
Can we compute the predictive probability function of ${\color{red}y_{i,n_i+1}}$ given $(y_{i,1}, \ldots, y_{i,n_i})$? \pause YES!

\end{frame}

\begin{frame}
\frametitle{Poisson-gamma GLMM - prediction}

By using the conditioning argument, we have
{\small
\begin{align*}
& f({\color{red}y_{i,n_i+1}}|y_{i,1}, \ldots, y_{i,n_i}) = \frac{h(y_{i,1}, \ldots, y_{i,n_i}, {\color{red}y_{i,n_i+1}})}{h(y_{i,1}, \ldots, y_{i,n_i})} \\
&= \left[ \frac{(\nu_{i,n_i+1})^{y_{i,n_i+1}}}{y_{i,n_i+1}!} \right] \times \frac{(\nu_{i,\bullet}+r)^{y_{i,\bullet}+r}}{\Gamma(y_{i,\bullet}+r)} \times \frac{\Gamma(y_{i,\bullet} + {\color{red}y_{i,n_i+1}} + r)}{(\nu_{i,\bullet} + {\color{red}\nu_{i,n_i+1}} + r)^{y_{i,\bullet} + {\color{red}y_{i,n_i+1}} + r}} \\
&= \frac{\Gamma(y_{i,\bullet} + {\color{red}y_{i,n_i+1}} + r)}{y_{i,n_i+1}!\Gamma(y_{i,\bullet}+r)} \times \\
& \left( \frac{\nu_{i,\bullet}+r}{\nu_{i,\bullet} + {\color{red}\nu_{i,n_i+1}} + r} \right)^{y_{i,\bullet}+r} \left( \frac{{\color{red}\nu_{i,n_i+1}}}{\nu_{i,\bullet} + {\color{red}\nu_{i,n_i+1}} + r} \right)^{{\color{red}y_{i,n_i+1}}} \\
&= \frac{\Gamma(y_{i,\bullet} + {\color{red}y_{i,n_i+1}} + r)}{y_{i,n_i+1}!\Gamma(y_{i,\bullet}+r)} \times \\
& \left( \frac{1}{1 + {\color{red}\nu_{i,n_i+1}}/(\nu_{i,\bullet}+r)} \right)^{y_{i,\bullet}+r} \left( \frac{{\color{red}\nu_{i,n_i+1}}/(\nu_{i,\bullet}+r)}{1 + {\color{red}\nu_{i,n_i+1}}/(\nu_{i,\bullet}+r)} \right)^{{\color{red}y_{i,n_i+1}}}
\end{align*}
}

\end{frame}

\begin{frame}
\frametitle{Poisson-gamma GLMM - prediction}

Therefore, we can confirm that
$$
{\color{red}y_{i,n_i+1}} \, | \, y_{i,1}, \ldots, y_{i,n_i} \sim \mathcal{NB}\left( y_{i,\bullet} + r, \, \frac{{\color{red}\nu_{i,n_i+1}}}{\nu_{i,\bullet} + r} \right),
$$
so that
$$
\mathbb{E}[y_{i,n_i+1}|y_{i,1}, \ldots, y_{i,n_i}] = (y_{i,\bullet}+r) \cdot \frac{\nu_{i,n_i+1}}{\nu_{i,\bullet}+r} = \frac{y_{i,\bullet}+r}{\nu_{i,\bullet}+r} \cdot \nu_{i,n_i+1}.
$$

\end{frame}

\begin{frame}
\frametitle{Poisson-gamma GLMM - prediction}

We may interpret the above predictive premium as follows:
\begin{align*}
\mathbb{E}[y_{i,n_i+1}|y_{i,1}, \ldots, y_{i,n_i}] &= \frac{y_{i,\bullet}+r}{\nu_{i,\bullet}+r} \cdot \nu_{i,n_i+1} \\
& \onslide<2->{= \mathbb{E}[\Theta_i|y_{i,1}, \ldots, y_{i,n_i}] \cdot \mathbb{E}[y_{i,n_i+1}]} \\
\omit\rlap{\onslide<3->{$\displaystyle = \frac{\text{(actual past \# of claims)}}{\text{(expected past \# of claims)}} \times \text{(expected current \# of claims)}.$}}
\end{align*}

\onslide<3->{Therefore, the predictive premium from a Poisson-gamma GLMM also shares the spirit of \textbf{credibility} premium!}

\end{frame}

% ---------------------------------------------------------------------
\begin{frame}
\frametitle{Poisson GLM - prediction}

In the case of a Poisson GLM, we assume that all the responses $\{y_{i,t}\}$ are independent so that $h(y_{i,1}, \ldots, y_{i,n_i}) = \prod_{t=1}^{n_i} f(y_{i,t})$, and
\begin{align*}
f({\color{red}y_{i,n_i+1}}|y_{i,1}, \ldots, y_{i,n_i}) &= \frac{h(y_{i,1}, \ldots, y_{i,n_i}, {\color{red}y_{i,n_i+1}})}{h(y_{i,1}, \ldots, y_{i,n_i})} \\
& \onslide<2->{= \frac{\prod_{t=1}^{n_i+1} f(y_{i,t})}{\prod_{t=1}^{n_i} f(y_{i,t})} = f({\color{red}y_{i,n_i+1}}),}
\end{align*}
\onslide<2->{and subsequently
$$
\mathbb{E}[y_{i,n_i+1}|y_{i,1}, \ldots, y_{i,n_i}] = \mathbb{E}[y_{i,n_i+1}] = \nu_{i,n_i+1}.
$$
It implies that in the case of Poisson GLM, past observations \textbf{DO NOT} affect the future premium due to the (implicitly assumed) independence among the observed claims.}

\end{frame}

% ---------------------------------------------------------------------
\begin{frame}
\frametitle{Out-of-sample validation for claim frequency}

It is natural to use the following measures, root-mean-squared error (RMSE) and Poisson deviance (PDEV) to assess the prediction performance:
\begin{align*}
\text{RMSE} : & \sqrt{\frac{1}{m} \sum_{i=1}^m (y_{i,T_i+1} - \hat{y}_{i,T_i+1})^2}, \\
\text{PDEV} : & \frac{2}{m} \sum_{i=1}^m \left( y_{i,T_i+1} \log \frac{y_{i,T_i+1}}{\hat{y}_{i,T_i+1}} - y_{i,T_i+1} + \hat{y}_{i,T_i+1} \right).
\end{align*}
We prefer a model with lower RMSE/PDEV.

\end{frame}

\begin{frame}
\frametitle{Out-of-sample validation for claim frequency}

\begin{table}[ht]
\centering
\begin{tabular}{lrr}
\toprule
 & Poisson & MVNB \\
\midrule
RMSE & 0.3565 & 0.3557 \\
PDEV & 0.4595 & 0.4564 \\
\hline\hline
\end{tabular}
\end{table}
It is shown that usage of Poisson-gamma GLMM slightly improves the prediction performance.

\end{frame}


%======================================================================
% Part 2
%======================================================================
\section{Credibility Premium with Static Random Effects}


\begin{frame}
\centering
\Large\textbf{Part 2: Credibility premium of $S$\\ with the static random effects}
\end{frame}

\begin{frame}{P\&C Insurance Claim Data Structure}
\begin{itemize}
	\item Let $N$ be the number of accidents and $Y_1, \cdots, Y_N$ be the size of claim per accident.
	\item For ratemaking in P\&C, we have to predict the cost of claims $S_{it}=\sum\limits_{k=1}^{N_{it}} Y_{itk}$.
	\item Policyholder $i$ is followed over time $t=1,\ldots,T_i$ years.
	\item Unit of analysis ``\textit{it}'' -- an insured driver $i$ over time $t$ (year)
	\item For each ``\textit{it}'', could have several claims, $k=0,1,\ldots,N_{it}$
	\item Have available information on: number of claims $N_{it}$, amount of claim $Y_{itk}$, exposure $e_{it}$ and covariates (explanatory variables) $x_{it}$
	\begin{itemize}
		\item covariates often include age, gender, vehicle type, building type, building location, driving history and so forth
	\end{itemize}
\end{itemize}
\end{frame}

\begin{frame}{The Two-Part Frequency-Severity Model}
\begin{itemize}
	\item We will model the pair $(N_{it}, C_{it})$ where
	$$
	C_{it} = \begin{cases}
		\dfrac{1}{N_{it}} \sum\limits_{k=1}^{N_{it}} Y_{itk}, & N_{it} > 0 \\
		\textit{Undefined}, & N_{it} = 0
	\end{cases}
	$$
	is the observed average claim size and $S_{it}=\sum\limits_{k=1}^{N_{it}} Y_{itk}$ is the observed {\color{purple} aggregate claim size}.
	\item The joint density of the {\color{purple} number of claims} and the {\color{purple} average claim size} can be decomposed as
	\begin{eqnarray*}
		f(N,C|\mathbf{x}) &=& f(N|\mathbf{x}) \times f(C|N,\mathbf{x}) \\
		\text{joint} &=& \text{frequency} \ \times \ \text{average severity}.
	\end{eqnarray*}
\end{itemize}
\end{frame}

\begin{frame}{Ratemaking with Predictive Analytics}
\vfill
$$
\begin{array}{ll}
	\text{[Observed]} & \{x_{it}, N_{it}, C_{it}, e_{it} | i=1,2,\ldots,I, t=1,2,\ldots,T_i\} \\
	& \qquad\quad\ \ \Bigg\Downarrow \ \text{\color{purple}Calibrated Model} \ \Bigg\Downarrow \\
	\text{[Predicted]} & \qquad\quad {\color{purple} N_{i,T_i+1},\, C_{i,T_i+1},\, \text{and } S_{i,T_i+1}}
\end{array}
$$
\pause
From now on, we drop the subscript $i$ for succinct expressions.
\vfill
\vfill
\end{frame}

\begin{frame}{Why Random Effects?}

The following hypothetical example shows us that we might capture the
\color{purple} unobserved heterogeneity in risk \color{black} by
observing the residual after controlling for the effects of observed
covariates, which can be explained in terms of \color{purple} random
effects \color{black} for policyholders A and B.

\begin{table}[h!t!]
\centering
\resizebox{\textwidth}{!}{
\begin{tabular}{c|c|c|c|c|c|c|c}
\multirow{2}{*}{Year} & \multirow{2}{*}{Gender} & \multirow{2}{*}{Age} & \multirow{2}{*}{Vehicle Size} & \multicolumn{2}{c|}{Policyholder A} & \multicolumn{2}{c}{Policyholder B} \\ \cline{5-8}
                      &                         &                      &                               & \# of Claim      & Amt of Claim     & \# of Claim      & Amt of Claim     \\ \hline
2015                  & M                       & 45                   & Medium                        & 0                & 0                & 1                & 500              \\
2016                  & M                       & 46                   & Medium                        & 0                & 0                & 2                & 4000             \\
2017                  & M                       & 47                   & Large                         & 1                & 200              & 1                & 8000             \\
\end{tabular}}
\end{table}
Note that random effects model has a good connection with  \color{purple} credibility theory  \color{black} as well.
\end{frame}

\begin{frame}{Frequency Model Specifications}
In actuarial practice, Poisson distribution has been used for the calibration of the number of claim with the presence of covariates as follows:
\begin{equation}\label{p2:eq:pois}
N_t|\mathbf{x}_t, e_t \overset{indep}{\sim} \mathcal{P}(\nu_t) \quad \text{where} \ \ \nu_t = e_t \exp(\mathbf{x}_t \alpha).
\end{equation}
Due to the longitudinal property of usual claim dataset, we may consider the following extension by incorporating random effects as follows:
\begin{equation}\label{p2:eq:poisre}
N_t|\theta^N \overset{indep}{\sim} \mathcal{P}(\nu_t \theta^N) \quad \text{where} \ \ \nu_t = e_t \exp(\mathbf{x}_t \alpha), \ \theta^N \sim \pi_N(\theta).
\end{equation}
\end{frame}

\begin{frame}{Frequency Model Specifications}
Let us consider the following frequency models:
\begin{itemize}
	\item For the count of claims (frequency)
	\setbeamertemplate{enumerate items}[default]
	\begin{enumerate}[(1)]
		\item Simple Poisson GLM: $N_t \sim \text{Pois}(e^{\mathbf{x}_t \alpha})$ so that $\E{N_t|\mathbf{x}_t} = e^{\mathbf{x}_t\alpha}$.
		\item {\color{blue} Poisson/Gamma Random Effect Model (= Multivariate NB Model)}
	\end{enumerate}
\end{itemize}
\end{frame}

\begin{frame}{Poisson/Gamma Random Effect Model}
Data likelihood: $N_t|\theta^N \overset{indep}{\sim} \mathcal{P}(\nu_t \theta^N)$ \\
Prior density: $\pi_N(\theta) \propto \theta^{r-1}e^{-\theta r}$ so that $\theta^N \sim \mathcal{G}(r,r)$ and $\E{\theta^N}=1, \Var{\theta^N}=\frac{1}{r}$. \\
Posterior density: \\
$\theta^N|n_1, n_2, \ldots, n_T \sim \mathcal{G}(\sum_{t=1}^T n_t + r, \sum_{t=1}^T \nu_t + r)$ \\
{\color{purple} Predictive distribution:}
\begin{equation}\label{p2:eq:mvnb}
N_{T+1}|n_1, n_2, \ldots, n_T \sim \mathcal{NB}\left(\sum_{t=1}^T n_t + r, \frac{\nu_{T+1}}{\sum_{t=1}^T \nu_t + r}\right)
\end{equation}
so that $\E{N_{T+1}|n_1, n_2, \ldots, n_T} = \frac{\sum_{t=1}^T n_t + r}{\sum_{t=1}^T \nu_t + r}\nu_{T+1}$.
\end{frame}

\begin{frame}{Implication of the Frequency Credibility Premium}
\begin{itemize}
	\item We can see that
	\begin{align*}
		&\E{N_{T+1}|\mathcal{F}_T} = \frac{r+{\color{purple}\sum_{t=1}^T n_t}}{r+{\color{purple}\sum_{t=1}^T \exp(\mathbf{x}_t\alpha)}} \times \exp(\mathbf{x}_{T+1}\alpha) \\
		&= \frac{\text{(actual past \# of claims)}}{\text{(expected past \# of claims)}} \times \text{(expected current \# of claims)}.
	\end{align*}
	\item In this sense, use of random effects helps us to penalize/reward a policyholder based on his/her driving experience.
	\item What about $\E{S_{T+1}|\mathcal{F}_T}$?
\end{itemize}
\end{frame}

\begin{frame}{Static Credibility Premium for S when $N \perp C$}
Frangos and Vrontos (2001) derived the following formula for (static) credibility premium of $S$ assuming $N \perp C$.
\begin{theorem}
Suppose $(N_1, N_2, \ldots, N_T)$ follows MVNB distribution as defined in \eqref{p2:eq:mvnb}. Moreover, let us assume that $\E{C_t|n_t} = \theta^C e^{\mathbf{x}_t \beta} = \theta^C \tilde{\mu}_t$. Then credibility premium of $S_{T+1}$ given $n_1, \ldots, n_T, c_1, \ldots, c_T$ is given as follows:
\begin{align*}
\mathbb{E}[S_{T+1}|\mathcal{F}_T] &= \mathbb{E}[N_{T+1}|n_1, \ldots, n_T] \times \mathbb{E}[C_{T+1}|\mathcal{F}_T] \\
&= \text{\emph{Cred. Freq. mean}} \ \times \text{\emph{Cred. Sev. mean}}
\end{align*}
\end{theorem}
\pause
{\color{purple} What if $N \not\perp C$?}
\end{frame}

\begin{frame}{Cross-sectional Dependence between $N$ and $C$}
Inspired by Garrido et al. (2016), we assume a specific form of cross-sectional dependence between $N_t$ and $C_t$ as follows:
\begin{equation}\label{p2:eq:crossdep}
\begin{aligned}
&N_t|\theta^N \overset{indep}{\sim} \mathcal{P}(\nu_t \theta^N) \quad \text{where} \ \ \nu_t = e_t \exp(\mathbf{x}_t \alpha), \ \theta^N \sim \mathcal{G}(r,r). \\
&C_t|\theta^C, n_t \overset{indep}{\sim} \mathcal{G}\left(\frac{n_t}{\psi}, \frac{n_t}{\psi\theta^C\mu_t}\right) \quad \text{where} \ \ \mu_t = \tilde{\mu}_t e^{{\color{purple}\gamma n_t}} = \exp(\mathbf{x}_t\beta + {\color{purple}\gamma n_t}), \\
&\theta^C \sim \mathcal{IG}(k+1,k).
\end{aligned}
\end{equation}
Since $\mu_t = \exp(\mathbf{x}_t\beta + {\color{purple}\gamma n_t})$, \\
$\E{S_t|\theta^N, \theta^C} \neq \E{N_t|\theta^N}\E{C_t|\theta^C, n_t}$.
\end{frame}

\begin{frame}{Severity Model Specifications}
Let us consider the following severity models:
\begin{itemize}
	\item For the average size of claims (severity)
	\setbeamertemplate{enumerate items}[default]
	\begin{enumerate}[(1)]
		\item Simple Gamma GLM: $C_t|n_t \sim \mathcal{G}(\frac{n_t}{\psi}, \frac{n_t}{\psi\mu_t})$ so that $\E{C_t|n_t, \mathbf{x}_t} = \mu_t = \tilde{\mu}_t e^{\gamma n_t}$ and $\frac{\Var{C_t|n_t, \mathbf{x}_t}}{\E{C_t|n_t, \mathbf{x}_t}^2} = \frac{\psi}{n_t}$
		\item {\color{blue} Gamma/I-gamma Random Effect Model (= Multivariate GP Model)}
	\end{enumerate}
\end{itemize}
\end{frame}

\begin{frame}{Gamma/I-gamma Random Effects Model}
Data likelihood: $C_t|\theta^C \overset{indep}{\sim} \mathcal{G}\left(\frac{n_t}{\psi}, \frac{n_t}{\psi\theta^C\mu_t}\right)$ \\
Prior density: $\pi_C(\theta) \propto \theta^{-k-2}e^{-k/\theta}$ so that $\theta^C \sim \mathcal{IG}(k+1,k)$ and $\E{\theta^C}=1, \Var{\theta^C}=\frac{1}{k-1}$. \\
Posterior density: \\
$\theta^C|\mathcal{F}_T \sim \mathcal{IG}(k+\sum_{t=1}^T n_t/\psi + 1, k+\sum_{t=1}^T \frac{s_t}{\psi\mu_t})$. \\
Note that $s_t = n_t \cdot c_t$ by definition.
\end{frame}

\begin{frame}{Gamma/I-gamma Random Effects Model}
\textcolor{purple}{Predictive distribution:}
\begin{align*}
&C_{T+1}|\mathcal{F}_T, n_{T+1} \\
&\sim \mathcal{GP}\left(k + \sum_{t=1}^T n_t/\psi + 1, \left[k + \sum_{t=1}^T \frac{s_t}{\psi\mu_t}\right]\frac{\psi\mu_{T+1}}{n_{T+1}}, \frac{n_{T+1}}{\psi}\right),
\end{align*}
with predictive mean
$$
\E{C_{T+1}|n_{T+1}, \mathcal{F}_T} = \frac{\left[k\psi + \sum_{t=1}^T s_t/\mu_t\right]}{k\psi + \sum_{t=1}^T n_t}\mu_{T+1}.
$$
\end{frame}

\begin{frame}{GP and GB2 Distributions}
\begin{itemize}
	\item $Y \sim \mathcal{GP}(a, \xi, \tau)$
	$$
	\Longrightarrow \quad f(y|a,\xi,\tau) = \frac{\Gamma(a+\tau)}{\Gamma(a)\Gamma(\tau)}\frac{\xi^a y^{\tau-1}}{(y+\xi)^{a+\tau}}
	$$
	\item $Y \sim \mathcal{GB}2(a, \xi, \tau, p)$
	$$
	\Longrightarrow \quad f(y|a,\xi,\tau,p) = \frac{\Gamma(a+\tau)}{\Gamma(a)\Gamma(\tau)}|p|\frac{\xi^{ap} y^{\tau p-1}}{(y^p+\xi^p)^{a+\tau}}
	$$
	\item Note that GP distribution is a special case of GB2 distribution where $p=1$. (Loss Models, Klugman et al. 2012)
\end{itemize}
\end{frame}

\begin{frame}{Static Credibility Premium for S when $N \not\perp C$}
\begin{theorem}
Suppose $(N_1, N_2, \ldots, N_T)$ follows MVNB distribution as defined in \eqref{p2:eq:mvnb}. Moreover, let us assume that \\
$\E{C_t|n_t} = \theta^C e^{\mathbf{x}_t\beta}\, {\color{purple}e^{n_t\gamma}} = \theta^C \tilde{\mu}_t\, {\color{purple}e^{n_t\gamma}}$. Then, the credibility premium of $S_{T+1}$ is given as follows:
\begin{equation}\label{p2:eq:credS}
\begin{aligned}
\E{S_{T+1}|\mathcal{F}_T} &= \E{\tilde{\mu}_{T+1}\theta^C|\mathcal{F}_T} \times \E{\nu_{T+1}\theta^N|\underline{n_T}} \\
&\quad \times e^\gamma\left[1-\left(\frac{\nu_{T+1}}{\sum_{t=1}^T \nu_t + r}\right)(e^\gamma - 1)\right]^{-\sum_{t=1}^T n_t - r - 1}.
\end{aligned}
\end{equation}
If $\gamma=0$, then \eqref{p2:eq:credS} is reduced as in Frangos and Vrontos (2001):
$$
\E{S_{T+1}|\mathcal{F}_T} = \E{\tilde{\mu}_{T+1}\theta^C|\mathcal{F}_T} \times \E{\nu_{T+1}\theta^N|\underline{n_T}}.
$$
\end{theorem}
\end{frame}

\begin{frame}{Static Credibility Premium for S when $N \not\perp C$}
\begin{itemize}
	\item According to the theorem, it can be observed that the \textbf{static} credibility premium of $S_{T+1}$ can be expressed as the product of
	\begin{itemize}
		\item {\color{purple}\textbf{static} credibility premium of $N_{T+1}$,}
		\item {\color{purple}\textbf{static} credibility premium of $C_{T+1}$,} and
		\item {\color{purple}adjustment factor which accounts for the dependence between frequency and the average severity.}
	\end{itemize}
	\item Under the MVNB distribution defined in \eqref{p2:eq:mvnb}, we have the following:
	$$
	\E{\nu_{T+1}\theta^N|\underline{n_T}} = \frac{r+\sum_{t=1}^T n_t}{r+\sum_{t=1}^T \nu_t}\nu_{T+1}.
	$$
	\item For convenience, let us denote \\
	$D_N(\gamma) = e^\gamma\left[1-\left(\frac{\nu_{T+1}}{\sum_{t=1}^T \nu_t + r}\right)(e^\gamma - 1)\right]^{-\sum_{t=1}^T n_t - r - 1}$ so that $D_N(0)=1$.
\end{itemize}
\end{frame}

\begin{frame}{Credibility Premium for the Compound Sum when $N \not\perp C$}
\begin{theorem}
For $D_N(\gamma)$ in \eqref{p2:eq:credS}, the following are true:
\setbeamertemplate{enumerate items}[default]
\begin{enumerate}[(i)]
	\item $D_N(\gamma)$ is well-defined if and only if $\gamma \leq \log\left(1+\frac{r+\sum_{t=1}^T \nu_t}{\nu_{T+1}}\right)$.
	\item $D_N(\gamma)$ is a strictly increasing function of $\gamma$.
	\item
	\small
	$$
	D_N(\gamma) = e^\gamma\left[1-\left(\frac{\nu_{T+1}}{r+\sum_{t=1}^T \nu_t}\right)(e^\gamma - 1)\right]^{-r-\sum_{t=1}^T n_t - 1}
	\begin{cases}
		=1 & \text{\emph{if}} \quad \gamma = 0 \\
		>1 & \text{\emph{if}} \quad \gamma > 0 \\
		<1 & \text{\emph{if}} \quad \gamma < 0
	\end{cases}
	$$
\end{enumerate}
\end{theorem}
In this regard, $D_N(\gamma)$ can be interpreted as an {\color{purple} informative measure of dependence} which considers not only the sign of dependence but also the magnitude of dependence.
\end{frame}

%======================================================================
% Part 3
%======================================================================
\section{Credibility Premium with Dynamic Random Effects}


\begin{frame}
\centering
\Large\textbf{Part 3: Credibility premium of $S$ with the dynamic random effects}
\end{frame}

%----------------------------------------------------------------------
\begin{frame}{Static Frequency Credibility Premium}
\begin{itemize}
\item Recall that in Part 1, the frequency credibility premium is given as follows:
\begin{align*}
&\E{N_{T+1}|\mathcal{F}_T} = \frac{r+\textcolor{purple}{\sum_{t=1}^T n_t}}{r+\textcolor{purple}{\sum_{t=1}^T \exp(\mathbf{x}_t\alpha)}}\times \exp(\mathbf{x}_{T+1}\alpha) \\
&= \frac{\text{(actual past \# of claims)}}{\text{(expected past \# of claims)}} \times \text{(expected current \# of claims)}.
\end{align*}
\pause
\item In this sense, using \textbf{static} random effects is somewhat equivalent to use
\textcolor{purple}{aggregated} actual/expected \# of claims for adjustments.
\pause
\item However, it means that there is \textcolor{purple}{no distinction on the age of claims.}
(i.e., $n_1=1$ or $n_5=1$ ends up with the same impacts in the credibility premium.)
\end{itemize}
\end{frame}

%----------------------------------------------------------------------
\begin{frame}{Disclaimer: STATIC Random Effects Models}
\begin{itemize}
\item In the formulation, we assumed that the random effect $\theta_i$ is
\textcolor{purple}{random but STATIC}. In other words, it is hard to capture possible evolution
on the unobserved heterogeneity in risks over time with such model.
\pause
\item This assumption leads to the predictive premium formula which assigns
\textcolor{purple}{the same weight for every $n_t$, $t=1,\ldots,T$}.
\pause
\item One can suggest \textcolor{purple}{DYNAMIC random effects} compound risk models (CRMs) to capture
\textcolor{purple}{possible evolution on the unobserved heterogeneity}, which follows in the later part of this presentation.
\end{itemize}
\end{frame}

%----------------------------------------------------------------------
\begin{frame}{Key Idea of Static Random Effects CRMs}
$$
\begin{array}{ccccccc}
\cdots & = & \theta^N & = & \theta^N & = & \theta^N \quad = \\
\textcolor{teal}{\downarrow} & & \textcolor{teal}{\downarrow} & & \textcolor{teal}{\downarrow} & & \textcolor{teal}{\downarrow} \\
\cdots & & N_{t-1} & & N_t & & N_{t+1} \\
 & & \textcolor{olive}{\Downarrow} & & \textcolor{olive}{\Downarrow} & & \textcolor{olive}{\Downarrow} \\
\cdots & & C_{t-1} & & C_t & & C_{t+1} \\
\textcolor{teal}{\uparrow} & & \textcolor{teal}{\uparrow} & & \textcolor{teal}{\uparrow} & & \textcolor{teal}{\uparrow} \\
\cdots & = & \theta^C & = & \theta^C & = & \theta^C \quad =
\end{array}
$$
\pause
\begin{itemize}
\item \textcolor{teal}{$\theta^N_{\phantom{t-1}}$ determines the conditional distribution of $N_{t-1}$ (respectively, C).}
\pause
\item \phantom{\parbox[t]{\linewidth}{$\theta^N_{t-1}$ and $N_{t-1}$ affect the filtering distribution of $\theta^N_t$ (respectively, C).}}
\pause
\item \textcolor{olive}{$N_t$ is assumed to have cross-sectional dependence with $C_t$.}
\end{itemize}
\end{frame}

%----------------------------------------------------------------------
\begin{frame}{Key Idea of Dynamic Random Effects CRMs}
$$
\begin{array}{cccccccc}
\cdots & \textcolor{violet}{\longrightarrow} & \theta^N_{\textcolor{blue}{t-1}} & \textcolor{violet}{\longrightarrow} & \theta^N_{\textcolor{blue}{t}} & \textcolor{violet}{\longrightarrow} & \theta^N_{\textcolor{blue}{t+1}} & \textcolor{violet}{\longrightarrow} \\
\textcolor{teal}{\downarrow} & \textcolor{violet}{\nearrow} & \textcolor{teal}{\downarrow} & \textcolor{violet}{\nearrow} & \textcolor{teal}{\downarrow} & \textcolor{violet}{\nearrow} & \textcolor{teal}{\downarrow} & \\
\cdots & & N_{t-1} & & N_t & & N_{t+1} & \\
 & & \textcolor{olive}{\Downarrow} & & \textcolor{olive}{\Downarrow} & & \textcolor{olive}{\Downarrow} & \\
\cdots & & C_{t-1} & & C_t & & C_{t+1} & \\
\textcolor{teal}{\uparrow} & \textcolor{violet}{\searrow} & \textcolor{teal}{\uparrow} & \textcolor{violet}{\searrow} & \textcolor{teal}{\uparrow} & \textcolor{violet}{\searrow} & \textcolor{teal}{\uparrow} & \\
\cdots & \textcolor{violet}{\longrightarrow} & \theta^C_{\textcolor{blue}{t-1}} & \textcolor{violet}{\longrightarrow} & \theta^C_{\textcolor{blue}{t}} & \textcolor{violet}{\longrightarrow} & \theta^C_{\textcolor{blue}{t+1}} & \textcolor{violet}{\longrightarrow}
\end{array}
$$
\pause
\begin{itemize}
\item \textcolor{teal}{$\theta^N_{\textcolor{blue}{t-1}}$ determines the conditional distribution of $N_{t-1}$ (respectively, C).}
\pause
\item \textcolor{violet}{$\theta^N_{\textcolor{blue}{t-1}}$ and $N_{t-1}$ affect the filtering distribution of $\theta^N_{\textcolor{blue}{t}}$ (respectively, C).}
\pause
\item \textcolor{olive}{$N_t$ is assumed to have cross-sectional dependence with $C_t$.}
\end{itemize}
\end{frame}

%----------------------------------------------------------------------
\begin{frame}{Building Blocks for Dynamic Random Effects CRMs}
\begin{itemize}
\item<2-> Innovation of the dynamic (gamma) random effects (Lukacs, 1955),
\item<3-> Frequency dynamic random effects models (Harvey and Fernandez, 1989),
\item<4-> Severity dynamic random effects models (Smith and Miller, 1986), and
\item<5-> Cross-sectional dependence between the frequency and severity (Garrido et al., 2016; Jeong and Valdez, 2020).
\end{itemize}
\end{frame}

%----------------------------------------------------------------------
\begin{frame}{Innovation for the Dynamic Random Effects}
\begin{lemma}[Lukacs, 1955]
Given two independent random variables $\theta_t \sim \mathcal{G}(u,v)$ and $Z_t \sim \mathcal{B}(wu,(1-w)u)$,
where $u>0, v>0, w \in (0,1)$ are constants, we have
$$
\theta_{t+1} := \frac{Z_t\theta_t}{w} \sim \mathcal{G}(wu, wv).
$$
\end{lemma}
\pause
\begin{itemize}
\item $\theta_t$: the dynamic random effect at time $t$,
\pause
\item $Z_t$: source of (random) innovation at time $t$ for the random effect.
\pause
\item It is easy to check that $\text{Cor}(\theta_t, \theta_{t+1}) = w$.
\end{itemize}
\pause
Indeed, if $w=1$, then $Z_t$ is almost surely 1 so that $\theta_{t+1}=\theta_t$, which is equivalent to the \textbf{static} random effects case.
\end{frame}

%----------------------------------------------------------------------
\begin{frame}{Frequency Dynamic Random Effects (HF, 1989)}
The joint dynamics of the count (i.e. frequency) $N_t$ and the random effect $\theta^N_t$ is as follows:
\setbeamertemplate{enumerate items}[default]
\begin{enumerate}[i.]
\item (Updating rule of $\theta^N_t$) Assume that at time $t>1$, the filtering distribution of $\theta^N_{t-1}$
given the past observations\footnote{We similarly define
$$
\underline{A_t} := (A_1,\cdots,A_t), \quad \underline{\theta^N_t} := (\theta^N_1,\cdots,\theta^N_t) \quad \text{and} \quad \underline{\nu_t} := (\nu_1,\cdots,\nu_t).
$$}
$\underline{n_{t-1}} := (n_1,\cdots,n_{t-1})$ is given by
$$
\theta^N_{t-1}|\underline{n_{t-1}} \sim \mathcal{G}\left(r_{t-1}, r'_{t-1}\right),
$$
where $r_1 = r'_1 = r$ and $r_{t-1}$, $r'_{t-1}$ are \textcolor{purple}{positive, deterministic functions of past counts and exogenous covariates observed up to time $t-1$}.
\end{enumerate}
\end{frame}

\begin{frame}{Frequency Dynamic Random Effects (HF, 1989)}
Then the new state variable $\theta^N_t$ is defined by
$$
\theta^N_t := \frac{\theta^N_{t-1}A_t}{p},
$$
where $p \in (0,1)$ is a predetermined constant and $A_t$ has the following predictive distribution:
$$
A_t \,|\, \underline{A_{t-1}}, \underline{n_{t-1}}, \underline{\theta^N_{t-1}} \sim \mathcal{B}\left(pr_{t-1}, (1-p)r_{t-1}\right).
$$
Note that the conditional distribution of $\theta^N_t$ given $\underline{n_{t-1}}$, called the predictive distribution, is given by:
\begin{equation}\label{p3:freqpred}
\theta^N_t|\underline{n_{t-1}} \sim \mathcal{G}\left(pr_{t-1}, pr'_{t-1}\right).
\end{equation}
\end{frame}

\begin{frame}{Frequency Dynamic Random Effects (HF, 1989)}
\setbeamertemplate{enumerate items}[default]
\begin{enumerate}[i.]
\setcounter{enumi}{1}
\item (Observation at time $t$) Given $\underline{n_{t-1}}$ and $\underline{\theta^N_t}$, observation $N_t$ has the following conditional distribution:
$$
N_t|\underline{n_{t-1}}, \underline{\theta^N_t} \sim \mathcal{P}\left(\nu_t\theta^N_t\right),
$$
where $\nu_t$ is a \textcolor{purple}{positive, deterministic function of exogenous covariates observed at time $t$}.
Thus by the Poisson-gamma conjugacy, the filtering distribution of $\theta^N_t$ given $\underline{n_t}$ is given by:
\begin{equation}\label{p3:freqfilt}
\theta^N_t|\underline{n_t} \sim \mathcal{G}\left(r_t, r'_t\right),
\end{equation}
with
$$
r_t := pr_{t-1} + n_t > 0 \quad and \quad r'_t := pr'_{t-1} + \nu_t > 0.
$$
\end{enumerate}
\end{frame}

\begin{frame}{Frequency Dynamic Random Effects (HF, 1989)}
\begin{itemize}
\item An immediate consequence is the following relationships concerning the predictive mean and variance of the random effect:
$$
\resizebox{\linewidth}{!}{$\displaystyle
\E{\theta^N_t \,|\, \underline{n_{t-1}}} = \E{\theta^N_{t-1} \,|\, \underline{n_{t-1}}}, \Var{\theta^N_t \,|\, \underline{n_{t-1}}} = \frac{1}{p}\Var{\theta^N_{t-1} \,|\, \underline{n_{t-1}}}.$}
$$
\item Thus, updating $\theta^N_{t-1}$ to $\theta^N_t$ has the effect of preserving the mean (as well as the parametric form of the distribution) when passing from the filtering distribution to the predictive distribution.
\item It introduces extra ``uncertainty'', through the multiplicative ``innovation'' $\frac{A_t}{p}$, which is of unit mean.
This transition has also the effect of inflating the variance by a factor of $1/p$.
\end{itemize}
\end{frame}

\begin{frame}{Frequency Dynamic Random Effects (HF, 1989)}
One can check that $N_{T+1}\,|\,\underline{n_T}$ still follows a negative binomial distribution with the following predictive mean:
{\small
\begin{equation}\label{p3:freqmean}
\begin{gathered}
\E{N_{T+1}\,|\,\underline{n_T}} = \nu_{T+1}\frac{r_T}{r'_T}
= \nu_{T+1}\frac{\sum\limits_{t=1}^T p^{T-t}n_t + p^Tr}{\sum\limits_{t=1}^T p^{T-t}\nu_t + p^Tr}
= \nu_{T+1}\left[a_0 + \sum_{t=1}^T a_t \frac{n_t}{\nu_t}\right], \\
\text{with } a_0 := \frac{p^Tr}{\sum\limits_{t=1}^T p^{T-t}\nu_t + p^Tr} \quad \text{and} \quad
a_t := \frac{p^{T-t}\nu_t}{\sum\limits_{t=1}^T p^{T-t}\nu_t + p^Tr}.
\end{gathered}
\end{equation}
}
\pause
\textcolor{purple}{What if $p=1$?}
\end{frame}

%----------------------------------------------------------------------
\begin{frame}{Severity Dynamic Random Effects (SM, 1986)}
The joint dynamics of the average claim amounts (i.e. severity) $C_t$ and the random effect $\theta^C_t$ is:
\setbeamertemplate{enumerate items}[default]
\begin{enumerate}[i.]
\item (Updating rule of $\theta^C_t$) Assume that at time $t>1$, the filtering distribution of $\theta^C_{t-1}$
given the past observations\footnote{We similarly define
$$
\underline{B_t} := (B_1,\cdots,B_t), \quad \underline{\theta^C_t} := (\theta^C_1,\cdots,\theta^C_t) \quad \text{and} \quad \underline{\mu_t} := (\mu_1,\cdots,\mu_t).
$$}
$\underline{c_{t-1}} := (c_1,\cdots,c_{t-1})$ and $\underline{n_{t-1}}$ is given by
$$
\theta^C_{t-1}|\underline{c_{t-1}}, \underline{n_{t-1}} \sim \mathcal{IG}\left(k_{t-1}, k'_{t-1}\right),
$$
where $k_1 - 1 = k'_1 = k$ and $k_{t-1}$, $k'_{t-1}$ are \textcolor{purple}{positive, deterministic functions of past counts and exogenous covariates observed up to time $t-1$}.
\end{enumerate}
\end{frame}

\begin{frame}{Severity Dynamic Random Effects (SM, 1986)}
Then the new state variable $\theta^C_t$ is defined by
$$
\theta^C_t := \frac{\theta^C_{t-1}q^*_t}{B_t},
$$
where $q^*_t = \frac{q(k_{t-1}-2)+1}{k_{t-1}-1} \in (0,1)$ and $B_t$ has the following predictive distribution:
$$
B_t \,|\, \underline{B_{t-1}}, \underline{c_{t-1}}, \underline{n_{t-1}}, \underline{\theta^C_{t-1}} \sim \mathcal{B}\left(q_tk_{t-1}, (1-q_t)k_{t-1}\right),
$$
with $q_t = \frac{q(k_{t-1}-2)+2}{k_{t-1}} \in (2/k_{t-1}, 1)$.\\
Note that the conditional distribution of $\theta^C_t$ given $\underline{c_{t-1}}$, called the predictive distribution, is given by:
\begin{equation}\label{p3:sevpred}
\theta^C_t|\underline{c_{t-1}}, \underline{n_{t-1}} \sim \mathcal{IG}\left(q_tk_{t-1}, q^*_tk'_{t-1}\right).
\end{equation}
\end{frame}

\begin{frame}{Severity Dynamic Random Effects (SM, 1986)}
\setbeamertemplate{enumerate items}[default]
\begin{enumerate}[i.]
\setcounter{enumi}{1}
\item (Observation at time $t$) Given $\underline{c_{t-1}}$, $\underline{n_t}$, and $\underline{\theta^C_t}$, observation $C_t$ has the following conditional distribution:
$$
C_t|\underline{c_{t-1}}, \underline{n_t}, \underline{\theta^C_t} \sim \mathcal{G}\left(\frac{n_t}{\psi}, \frac{n_t}{\psi\mu_t\theta^C_t}\right),
$$
where $\mu_t$ is a \textcolor{purple}{positive, deterministic function of exogenous covariates observed at time $t$}.
Thus by the gamma-IG conjugacy, the filtering distribution of $\theta^C_t$ given $\underline{c_t}$ and $\underline{n_t}$ is given by:
\begin{equation}\label{p3:sevfilt}
\theta^C_t|\underline{c_t}, \underline{n_t} \sim \mathcal{IG}\left(k_t, k'_t\right),
\end{equation}
with
$$
k_t := q_tk_{t-1} + n_t/\psi > 0 \quad and \quad k'_t := q^*_tk'_{t-1} + \frac{s_t}{\psi\mu_t} > 0.
$$
\end{enumerate}
\end{frame}

\begin{frame}{Severity Dynamic Random Effects (SM, 1986)}
\begin{itemize}
\item An immediate consequence is the following relationships concerning the predictive mean and variance of the random effect:
$$
\resizebox{\linewidth}{!}{$\displaystyle
\E{\theta^C_t \,|\, \mathcal{F}_{t-1}} = \E{\theta^C_{t-1} \,|\, \mathcal{F}_{t-1}}, \Var{\theta^C_t \,|\, \mathcal{F}_{t-1}} = \frac{1}{q}\Var{\theta^C_{t-1} \,|\, \mathcal{F}_{t-1}}.$}
$$
\item Thus, updating $\theta^C_{t-1}$ to $\theta^C_t$ has the effect of preserving the mean (as well as the parametric form of the distribution) when passing from the filtering distribution to the predictive distribution.
\item It introduces extra ``uncertainty'', through the multiplicative ``innovation'' $\frac{B_t}{q}$, which is of unit mean.
This transition has also the effect of inflating the variance by a factor of $1/q$.
\end{itemize}
\end{frame}

\begin{frame}{Severity Dynamic Random Effects (SM, 1986)}
One can check that $C_{T+1}\,|\,\mathcal{F}_T, n_{T+1}$ still follows a GP distribution with the following predictive mean:
\begin{equation}\label{p3:sevmean}
\begin{gathered}
\E{C_{T+1}\,|\,\mathcal{F}_T, n_{T+1}} = \mu_{T+1}\frac{k'_T}{k_T-1}
= \mu_{T+1}\left[b_0 + \sum_{t=1}^T b_t\frac{s_t}{\mu_t}\right], \\
\text{with } b_0 := \frac{k\prod_{s=1}^T q^*_s}{k_T-1} \quad \text{and} \quad
b_t := \frac{\prod_{s=t+1}^T q^*_s/\psi}{k_T-1}.
\end{gathered}
\end{equation}
\pause
\textcolor{purple}{What if $q=1$? (Note that $q^*_s=q_s=1$ for all $s$ if $q=1$, by definition)}
\end{frame}

%----------------------------------------------------------------------
\begin{frame}{Dynamic Credibility Premium for S when $N \not\perp C$}
\begin{theorem}
Suppose the predictive means of $N_{T+1}$ and $C_{T+1}$ are given as in (\ref{p3:freqmean}) and (\ref{p3:sevmean}), respectively.
Then, the (dynamic) credibility premium of $S_{T+1}$ is given as follows:
\begin{equation}\label{p3:dynS}
\begin{split}
\E{S_{T+1}|\mathcal{F}_T} &= \tilde{\mu}_{T+1}\frac{k'_T}{k_T-1} \times \nu_{T+1}\frac{r_T}{r'_T} \\
&\quad \times e^\gamma\left[1-\left(\frac{\nu_{T+1}}{r'_T}\right)(e^\gamma-1)\right]^{-r_T-1}.
\end{split}
\end{equation}
If $\gamma=0$, then (\ref{p3:dynS}) is reduced as follows:
$$
\E{S_{T+1}|\mathcal{F}_T} = \tilde{\mu}_{T+1}\frac{k'_T}{k_T-1} \times \nu_{T+1}\frac{r_T}{r'_T}.
$$
\end{theorem}
\end{frame}

\begin{frame}{Dynamic Credibility Premium for S when $N \not\perp C$}
\begin{itemize}
\item According to the theorem, it can be observed that the \textbf{dynamic} credibility premium of $S_{T+1}$ can be expressed as the product of
\begin{itemize}
\item \textcolor{purple}{\textbf{dynamic} credibility premium of $N_{T+1}$},
\item \textcolor{purple}{\textbf{dynamic} credibility premium of $C_{T+1}$}, and
\item \textcolor{purple}{adjustment factor which accounts for the dependence between frequency and the average severity.}
\end{itemize}
\end{itemize}
\end{frame}

%----------------------------------------------------------------------
\begin{frame}{Data}
\rmfamily
We use Wisconsin Local Government Property Insurance Fund (LGPIF) dataset.
While claim information on multiple lines of business is available, we focus on IM (inland marine) claims and their associated policy characteristics from the selected dataset.

\vspace{1em}
\centering
\scriptsize
\begin{tabular}{l|lrrr}
\hline
\makecell[l]{Categorical\\variables} & Description & & \multicolumn{2}{c}{Proportions} \\
TypeCity & Indicator for city entity: & Y=1 & \multicolumn{2}{c}{14.51 \%} \\
TypeCounty & Indicator for county entity: & Y=1 & \multicolumn{2}{c}{5.92 \%} \\
TypeMisc & Indicator for miscellaneous entity: & Y=1 & \multicolumn{2}{c}{10.78 \%} \\
TypeSchool & Indicator for school entity: & Y=1 & \multicolumn{2}{c}{29.1 \%} \\
TypeTown & Indicator for town entity: & Y=1 & \multicolumn{2}{c}{16.6 \%} \\
TypeVillage & Indicator for village entity: & Y=1 & \multicolumn{2}{c}{23.09 \%} \\
\hline
\makecell[l]{Continuous\\variables} & & Minimum & Mean & Maximum \\
CoverageIM & Log coverage amount of IM claim & 0.00 & 0.87 & 46.75 \\
lnDeductIM & Log deductible amount of IM claim & 0.00 & 5.34 & 9.21 \\
\hline\hline
\end{tabular}
\end{frame}

%----------------------------------------------------------------------
\begin{frame}{Benchmark models}
For comparison, we introduce three benchmark models, which are special or limiting cases of the proposed dynamic random effects model:
\begin{itemize}
\item The naive GLMs: $p=q=1$, $r=k=\infty$, and $\eta=0$,
\item Frequency-severity dependent GLMs (Garrido et al., 2016): $p=q=1$, $r=k=\infty$,
\item Static random effects model (Jeong and Valdez, 2020): $p=q=1$, and
\item Dynamic random effects model.
\end{itemize}
\end{frame}

\begin{frame}{Benchmark models}
For estimation of the parameters, we take a two-step approach so that
\begin{itemize}
\item the fixed effects parameters are estimated first due to the common mean structure via usual GLMs, and the dependence structure is considered later.
In other words, we fit a usual Poisson/Gamma GLMs to estimate $\alpha$, $\beta$, and $\gamma$ first,
\item and the remaining parameters for the dependence structures ($p$, $q$, $r$, $k$, and $\psi$) are estimated afterward.
\end{itemize}
\end{frame}

%----------------------------------------------------------------------
\begin{frame}{Regression Estimates of the Fixed Effects}
\centering
\scriptsize
\begin{tabular}{lrrrrrr}
\toprule
 & \multicolumn{2}{c}{Frequency} & \multicolumn{4}{c}{Severity} \\
\cmidrule(lr){2-3}\cmidrule(lr){4-7}
 & & & \multicolumn{2}{c}{Independent} & \multicolumn{2}{c}{Dependent} \\
\cmidrule(lr){4-5}\cmidrule(lr){6-7}
 & Estimate & p-value & Estimate & p-value & Estimate & p-value \\
\midrule
(Intercept) & -4.2571 & 0.0000 & 9.4271 & 0.0000 & 10.1312 & 0.0000 \\
TypeCity & 0.9673 & 0.0000 & 1.0040 & 0.0339 & 1.3007 & 0.0012 \\
TypeCounty & 1.8747 & 0.0000 & 1.3164 & 0.0060 & 1.3152 & 0.0009 \\
TypeMisc & -2.7453 & 0.0068 & -1.3007 & 0.6018 & -1.2977 & 0.5252 \\
TypeSchool & -0.9174 & 0.0008 & -0.2046 & 0.7617 & -0.1191 & 0.8300 \\
TypeTown & -0.4129 & 0.1254 & 0.2571 & 0.6981 & 0.3588 & 0.5089 \\
CoverageIM & 0.0687 & 0.0000 & 0.0008 & 0.9712 & 0.0332 & 0.0746 \\
lnDeductIM & 0.1363 & 0.0030 & -0.1336 & 0.4481 & -0.1744 & 0.2278 \\
Count & & & & & -0.4538 & 0.0000 \\
\hline\hline
\end{tabular}
\end{frame}

%----------------------------------------------------------------------
\begin{frame}{Parameter Estimation for the Dependence Structure}
\begin{itemize}
\item Static random effects model: $\hat{r}=0.6143$, $\hat{k}=1.2615$, and $\hat{\psi}=1.5145$.
\item Dynamic random effects model: $\hat{p}=0.8202$, $\hat{q}=0.7615$, $\hat{r}=0.9228$, $\hat{k}=1.8761$, and $\hat{\psi}=1.4671$.
\end{itemize}
\end{frame}

\begin{frame}{Parameter Estimation for the Dependence Structure}
\begin{itemize}
\item Recall that
\begin{align*}
\theta^N_1 &\sim \mathcal{G}\left(r,r\right), & \E{\theta^N_1} &= 1,\ \Var{\theta^N_1} = \frac{1}{r}, \\
\theta^C_1 &\sim \mathcal{IG}\left(k+1,k\right), & \E{\theta^C_1} &= 1,\ \Var{\theta^C_1} = \frac{1}{k-1}.
\end{align*}
Since $r > k-1$ with the estimates from both static and dynamic models, it implies that the
\textcolor{purple}{variability of $\theta^C_1$}, the (a priori) \textcolor{purple}{unobserved heterogeneity for the severity} part,
is greater than that of $\theta^N_1$, the (a priori) unobserved heterogeneity for the \textcolor{purple}{frequency} part.
\item Secondly, $p$ is estimated to be greater than $q$ in the dynamic model, which implies the
\textcolor{purple}{decaying impact} according to the ages of claims in the determination of the credibility factors would be much
\textcolor{purple}{stronger in the severity} part.
\end{itemize}
\end{frame}

%----------------------------------------------------------------------
\begin{frame}{Out-of-sample Validation}
\centering
\small
\begin{tabular}{c|cccc}
\hline\hline
 & Naive & Garrido & Static & Dynamic \\
\hline
RMSE & 9272.96 & \textcolor{blue}{6433.21} & 6613.20 & 6608.58 \\
MAE & 1345.41 & 1241.95 & 1190.02 & \textcolor{blue}{1166.10} \\
Portfolio Mean & 979.94 & 905.95 & 884.05 & \textcolor{blue}{852.36} \\
\hline\hline
\multicolumn{5}{c}{* Note that the actual portfolio mean is 645.25.}
\end{tabular}
\end{frame}

%----------------------------------------------------------------------
\begin{frame}{Implication of Static vs. Dynamic Random Effects}
Let us consider a hypothetical scenario where
\begin{itemize}
\item $\nu_t = 0.2$ and $\mu_t = 15000$ for $t=1,\ldots,5$, and
\item $r=1$, $k=2$, $\gamma=0$, and $\psi=1.5$.
\end{itemize}
Suppose that there was \textcolor{purple}{exactly one claim} by a policyholder in one of the years 1 through 4
(i.e., $\sum_{t=1}^4 n_t = 1$). As shown in the following table, we observe decaying effects of past claims in the
calculation of credibility factors under the proposed dynamic model, whereas the static random effects model disregards such information.
\end{frame}

\begin{frame}{Implication of Static vs. Dynamic Random Effects}
\centering
\small
\begin{tabular}{ccccc}
\toprule
 & \multicolumn{2}{c}{Dynamic ($p=q=0.8$)} & \multicolumn{2}{c}{Static ($p=q=1$)} \\
\cmidrule(lr){2-3}\cmidrule(lr){4-5}
Claim year & Frequency & Severity & Frequency & Severity \\
\midrule
1 & 0.9216 & 1.0309 & 1.1111 & 1.0309 \\
2 & 1.0496 & 1.0347 & 1.1111 & 1.0309 \\
3 & 1.2096 & 1.0385 & 1.1111 & 1.0309 \\
4 & 1.4096 & 1.0421 & 1.1111 & 1.0309 \\
\hline\hline
\end{tabular}
\end{frame}

%----------------------------------------------------------------------
\begin{frame}{Concluding Remarks}
\begin{itemize}
\item Throughout the presentation, we explored ratemaking frameworks that consider \textcolor{purple}{longitudinality} as well as
\textcolor{purple}{dependence between frequency and severity} components.
\item In both cases (where the random effects are either static or dynamic), we could obtain
\textcolor{purple}{more general formula for credibility premium of the compound sum} with dependence between the frequency and severity.
\item Use of time-varying random effects processes enables us to account for \textcolor{purple}{the seniority of claims}
in the prediction formulas and lead to many nice features.
\end{itemize}
\end{frame}

\end{document}
